\documentclass[french]{book}
\usepackage{teach}
\usepackage{verbatim}
\usepackage{amsmath,amsfonts,amssymb}
\usepackage{pgfplots}
\pgfplotsset{compat=newest}
%\usepackage{geometry}
%\geometry{top=1cm, bottom=1cm, left=2cm, right=2cm}
\usepackage[utf8]{inputenc}
\usepackage[french]{babel}
\redefinecolor{thm}{title}{bgcolor}{alizarin}
\redefinecolor{prop}{title}{bgcolor}{meatbrown}
\redefinecolor{defn}{title}{bgcolor}{olivine}
\definecolor{alizarin}{rgb}{0.82, 0.1, 0.26}
\definecolor{meatbrown}{rgb}{0.9, 0.72, 0.23}
\definecolor{olivine}{rgb}{0.6, 0.73, 0.45}
\usepackage{pgf,tikz,tkz-tab}

\usepackage{mathrsfs}
\usetikzlibrary{arrows}

\begin{document}
 \section*{Devoir surveillé }
Soit $(O,\overrightarrow{i},\overrightarrow{j})$ un repère orthonormé. Toutes les coordonnées seront exprimées par rapport à ce repère. \\


1) Soit $\Delta$ une droite passant par deux points, $A=(1,2)$ et $B=(0,1)$. Déterminer l'équation cartésienne de la droite $\Delta$. \\


2) Soit $\Delta_1 $ une droite passant par $B$ et étant dirigé par le vecteur $\overrightarrow{u} \left( \begin{array}{c}
1 \\
3 \\
\end{array} \right). \\
$
Déterminer l'équation cartésienne de la droite $\Delta_1$.

3) Soit $\Delta_3 $ une droite ayant pour équation cartésienne $x-y+1=0$. Déterminer la pente de la droite  $\Delta_3 $. \\

4) Soit $\Delta_4$ une droite ayant pour équation cartésienne $-3x-y+1=0$. Déterminer l'équation réduite de cette droite et tracer la droite $\Delta_4$  dans le repère.\\

5) Soit un point $C=(1000,1001)$, dire si $A,B$ et $C$ sont alignés ou non. \\

6) Soit $\Delta_5$ une droite d'équation réduite $y=x+3$, dire si les deux droites $\Delta_5$ et $\Delta_3 $ sont strictement parallèles, confondus ou sécantes. \\

7) On suppose que les droites $\Delta_3$ et $\Delta_4$ sont sécantes. On cherche le point d'intersection $I$ de ces deux droites. Pour cela résoudre le système suivant et interprété le résultat: \\

$$
\left \{
\begin{array}{rcl}
x-y+1&=&0 \\
-3x-y+1&=&0
\end{array}
\right. 
$$

\definecolor{uuuuuu}{rgb}{0.26666666666666666,0.26666666666666666,0.26666666666666666}

\begin{center}
\begin{tikzpicture}[line cap=round,line join=round,>=triangle 45,x=1cm,y=1cm]
\begin{axis}[
x=1cm,y=1cm,
axis lines=middle,
ymajorgrids=true,
xmajorgrids=true,
xmin=-7.796395588853748,
xmax=8.608153986767181,
ymin=-4.83064566062787,
ymax=5.682847889614482,
xtick={-7,-6,...,8},
ytick={-4,-3,...,5},]
\clip(-7.796395588853748,-4.83064566062787) rectangle (8.608153986767181,5.682847889614482);
\draw [->,line width=0.5pt] (0,0) -- (1,0);
\draw [->,line width=0.5pt] (0,0) -- (0,1);
\begin{scriptsize}
\draw [fill=uuuuuu] (0,0) circle (0.7pt);
\draw[color=uuuuuu] (-0.19926640761063366,-0.2519556086633577) node {$O$};
\draw[color=black] (0.5298246846391854,-0.31549904668653216) node {$\overrightarrow{i}$};
\draw[color=black] (-0.3701027639206409,0.4448284421682512) node {$\overrightarrow{j}$};
\end{scriptsize}
\end{axis}
\end{tikzpicture}
\end{center}

\end{document}

